\scanpage{362}
\mspage{101}

d'un système inductif dans $\mathrm{Cat}$ (pas seulement dans $\mathrm{Cat\,hot}$)
\[
(C_n)_{n\in\mathbb N},
\]
d'où \unsure{on|en} récupère le ind-objet $\underline{\underline C}$ dans $\mathrm{Cat\,hot}$ en passant au syst. inductif correspondant dans $\mathrm{Cat\,hot}$, puis à l'ind-objet associé.

Pour construire le système inductif $C_n$, il suffit de prendre une chaine infinie
\diagram{\begin{tikzcd}[ampersand replacement=\&,column sep=.7em]
t_0\&t_1\arrow[l]\&t_2\arrow[l]\&\cdots\&t_n\arrow[l]\&t_{n+1}\arrow[l]\&\cdots\arrow[l]
\end{tikzcd}}
de morphismes dans $T$, telle que l'ens des $t_i$ soit cofinal \inserted{à gauche} dans $T_0$. \deleted{Ainsi} On va supposer de plus $C$ \uline{strictement fibrée} sur \deleted{$\unsure{\Sigma|I}$} $T$ (i.e., provenant d'un véritable foncteur $T^\circ\longrightarrow\mathrm{Cat}$,) i.e. on peut choisir les foncteurs images inverses \deleted{$\unsure{\varphi}$}
\[
\varphi^*:C_{t'}\mathrel{\unsure{-|\to}}C_t
\]
de façon transitive stricte, pas seulement transitive à iso. près) -- cas auquel on peut tjrs se ramener d'ailleurs par l'astuce de Giraud. Cette hypothèse est d'ailleurs vérifiée dans le cas $\underline{\mathrm{Ch}}_\Sigma(X;a,b)$ sur $\underline{\mathrm{Ty}}$. Ainsi, \deleted{\unsure{les}}
\[
(C_{t_i})_{i\in\mathbb N}
\]
forment un système \deleted{projectif} \inserted{inductif} \unsure{de|des} \unsure{catégorie|catégories}, \unsure{indexés|indexées|indexé} par $\mathbb N$, le morphisme de transition
\[
C_{\unsure{n}}\longrightarrow C_{n'}\qquad\text{i.e.}\quad C_{t_n}\mathrel{-}C_{t_{n'}}
\]

% archive sheet number: 361
